\chapter*{Referencias bibliográficas}
\addcontentsline{toc}{chapter}{Referencias bibliográficas}

\begin{thebibliography}{99}

\bibitem{Dashtdar2022}
M. Dashtdar et al., ``Frequency control of the islanded microgrid based on the model predictive control optimized by PSO,'' \textit{IET Renewable Power Generation}, vol. 16, no. 10, pp. 2139--2152, 2022.

\bibitem{Zishan2022}
F. Zishan et al., ``Efficient PID control design for frequency regulation in an independent microgrid based on the hybrid PSO-GSA algorithm,'' \textit{Electronics}, vol. 11, no. 23, p. 3886, 2022.

\bibitem{Gulzar2022}
M. M. Gulzar et al., ``Load frequency control strategies in renewable energy-based hybrid power systems: A review,'' \textit{Energies}, vol. 15, no. 10, p. 3488, 2022.

\bibitem{Kennedy1995}
J. Kennedy and R. Eberhart, ``Particle swarm optimization,'' in \textit{Proceedings of the International Conference on Neural Networks}, vol. 4, pp. 1942--1948, 1995.

\bibitem{Nepal2023}
T. Nepal et al., ``Models for hydropower plant: A review,'' in \textit{Proceedings of the 64th SIMS Conference on Simulation and Modelling}, 2023.

\bibitem{Alayi2021}
R. Alayi et al., ``Optimal load frequency control of island microgrids via a PID controller in the presence of wind and photovoltaic turbines,'' \textit{Sustainability}, vol. 13, no. 19, p. 10728, 2021.

\bibitem{AlBarazi2025}
A. A. Al Barazi and E. Aykut, ``Automated tuning of PID controllers using a generalized discrete PID structure with particle swarm optimization,'' \textit{IEEE Access}, vol. 13, pp. 185994--186013, 2025.

\bibitem{Joseph2022}
S. B. Joseph et al., ``Metaheuristic algorithms for PID controller parameters tuning: review, approaches and open problems,'' \textit{Heliyon}, vol. 8, no. 5, p. e09399, 2022.

\bibitem{Jasim2023}
A. M. Jasim et al., ``Optimal PID controller design for frequency regulation in independent microgrids using a hybrid PSO-GWO algorithm,'' \textit{Energies}, vol. 16, no. 15, p. 5748, 2023.

\bibitem{Rinaldi2021}
F. Rinaldi et al., ``Economic feasibility analysis and optimization of hybrid renewable energy systems for rural electrification in Peru,'' \textit{Clean Technologies and Environmental Policy}, vol. 23, no. 3, pp. 731--748, 2021.

\bibitem{Quispe2024}
J. C. Quispe, A. E. Obispo, y F. J. Alcántara, ``Economic feasibility assessment of microgrids with renewable energy sources in Peruvian rural areas,'' \textit{Clean Technologies and Environmental Policy}, vol. 26, no. 5, pp. 1415--1438, 2024.

\bibitem{Juanpera2021}
M. Juanpera et al., ``Renewable-based electrification for remote locations. Does short-term success endure over time? A case study in Peru,'' \textit{Renewable and Sustainable Energy Reviews}, vol. 146, p. 111177, 2021.

\end{thebibliography}
